\section{Multihead CNN model -- Lightweight OCR digits}

In this project, our system need to detect the timer display on Vietnamese traffic lights shows a one- to three-digit countdown
(0--999) and it must be recognised in real time on low-cost hardware such as the
Raspberry~Pi~5.  This section presents a compact, four-head convolutional neural network
(CNN) that is derived from the multi-digit recognition framework of
Goodfellow~et~al.~\cite{goodfellow2014multidigitnumberrecognitionstreet} and specifically re-engineered
for sub-millisecond CPU inference.

%-----------------------------------------------------------------------
\subsection{Probabilistic Formulation}
\label{sec:cnn_prob}

Following Goodfellow~et~al.~\cite{goodfellow2014multidigitnumberrecognitionstreet}, the recognition
task is cast as a structured probabilistic model.  Let~$\mathbf{X}$ denote the
input image and let the output be a digit sequence
$\mathbf{S} = (s_1, s_2, \dots, s_n)$ of length~$n$.
The sequence is modelled by a length variable~$L$ and a fixed number~$N$
of digit variables~$S_1, \dots, S_N$.  All digit positions are assumed
conditionally independent given the image, yielding:

\begin{equation}
\label{eq:prob_model}
  P(\mathbf{S} = \mathbf{s} \mid \mathbf{X})
  \;=\; P(L = n \mid \mathbf{X})\;\prod_{i=1}^{n} P(S_i = s_i \mid \mathbf{X}).
\end{equation}

In the original work~\cite{goodfellow2014multidigitnumberrecognitionstreet}, $N{=}5$ digit heads are
used for street numbers of up to five digits.  Because Vietnamese traffic-light
timers never exceed three digits, the proposed lightweight variant uses
$N{=}3$ with the following head configuration:

\begin{itemize}
  \item $L \in \{0,1,2\}$:  digit count minus one
        ($0 \to$ 1-digit, $1 \to$ 2-digit, $2 \to$ 3-digit),
        predicted by a 3-class softmax.
  \item $S_1 \in \{0,\dots,9,\text{BLANK}\}$:  hundreds digit (11 classes).
  \item $S_2 \in \{0,\dots,9,\text{BLANK}\}$:  tens digit    (11 classes).
  \item $S_3 \in \{0,\dots,9\}$:                 units digit  (10 classes,
        always present).
\end{itemize}

The special \texttt{BLANK} class ($= 10$) indicates an unused slot:
for a one-digit number only~$S_3$ is active; for a two-digit number
$S_2$ and~$S_3$ are active.

%-----------------------------------------------------------------------
\subsection{Network Architecture -- Edge-Optimised Design}
\label{sec:cnn_arch}

The original Goodfellow~et~al.\ architecture employs eight convolutional layers,
one locally connected layer, and two dense layers with 3\,072 units each,
totalling tens of millions of parameters and requiring DistBelief for training.
Such a model is impractical for real-time inference on a Raspberry~Pi~5 CPU.

The proposed \emph{lightweight} variant retains the multi-head output structure
but applies three key design transformations, each yielding a dramatic reduction in
multiply-accumulate operations (MACs) and parameter count while preserving accuracy
for the constrained 0--999 digit domain.

%--- Separable Convolutions -------------------------------------------------
\subsubsection{Depthwise Separable Convolutions}
\label{sec:separable}

Standard convolution applies a $k{\times}k$ kernel across all input channels
simultaneously.  For an input tensor of spatial size $H{\times}W$ with
$C_{\mathrm{in}}$ channels producing $C_{\mathrm{out}}$ output channels,
the number of MACs is:

\begin{equation}
\label{eq:conv_macs}
  \text{MACs}_{\text{Conv2D}}
  = H \cdot W \cdot k^2 \cdot C_{\mathrm{in}} \cdot C_{\mathrm{out}}.
\end{equation}

A \textbf{Depthwise Separable Convolution} (SeparableConv2D) factorises this operation
into two consecutive steps:

\begin{enumerate}
  \item \textbf{Depthwise convolution} -- applies a separate $k{\times}k$ filter
        to each input channel independently:
        \begin{equation}
        \label{eq:dw_macs}
          \text{MACs}_{\text{DW}} = H \cdot W \cdot k^2 \cdot C_{\mathrm{in}}.
        \end{equation}
  \item \textbf{Pointwise convolution} -- a $1{\times}1$ convolution that mixes
        the channel responses:
        \begin{equation}
        \label{eq:pw_macs}
          \text{MACs}_{\text{PW}} = H \cdot W \cdot C_{\mathrm{in}} \cdot C_{\mathrm{out}}.
        \end{equation}
\end{enumerate}

The total cost becomes:

\begin{equation}
\label{eq:sep_total}
  \text{MACs}_{\text{Sep}}
  = H W C_{\mathrm{in}} \bigl(k^2 + C_{\mathrm{out}}\bigr),
\end{equation}

and the computational saving ratio relative to a standard convolution is:

\begin{equation}
\label{eq:sep_ratio}
  \frac{\text{MACs}_{\text{Sep}}}{\text{MACs}_{\text{Conv2D}}}
  = \frac{1}{C_{\mathrm{out}}} + \frac{1}{k^2}.
\end{equation}

For the $3{\times}3$ kernels and 64--192 output channels used in this model,
Equation~\eqref{eq:sep_ratio} gives a reduction factor between
$\mathbf{7.8\times}$ and $\mathbf{8.5\times}$, enabling real-time execution
on an ARM Cortex-A76 CPU without a GPU.

In the proposed architecture, only the first block uses a standard
\texttt{Conv2D} (because the input has a single channel, where depthwise
separation offers no benefit); blocks~2--5 all use \texttt{SeparableConv2D}.

%--- GAP instead of Flatten ------------------------------------------------
\subsubsection{Global Average Pooling (GAP)}
\label{sec:gap}

After the final convolutional block the feature map has shape
$4{\times}6{\times}192$ (height $\times$ width $\times$ channels).
A na\"ive \texttt{Flatten} would produce a $4608$-dimensional vector; a
subsequent dense layer of width~1024 alone would add
$4608 \times 1024 \approx 4.7$~million parameters -- exceeding the entire
rest of the network.

Instead, {Global Average Pooling} (GAP) collapses each channel to its
spatial mean:

\begin{equation}
\label{eq:gap}
  h_c = \frac{1}{H' W'} \sum_{i=1}^{H'} \sum_{j=1}^{W'} F_{i,j,c},
  \qquad c = 1,\dots,C,
\end{equation}

where $F \in \mathbb{R}^{H' \times W' \times C}$ is the last feature map.
This yields a compact 192-dimensional descriptor with \emph{zero} learnable
parameters, and also acts as a structural regulariser that discourages
overfitting.

%--- Architecture Table ---------------------------------------------------
\subsubsection{Layer-by-Layer Specification}
\label{sec:arch_table}

Table~\ref{tab:arch} summarises the full architecture.  Input images are
grayscale, $64{\times}96$ pixels ($H{\times}W{\times}1$).

\begin{table}[htbp]
\centering
\caption{Layer-by-layer architecture of the lightweight multi-head CNN
         ($\approx 109$\,K parameters).  ``Sep'' denotes SeparableConv2D;
         ``BN'' is Batch Normalisation.}
\label{tab:arch}
\renewcommand{\arraystretch}{1.15}
\begin{tabular}{c l l r r}
\hline
\textbf{Block} & \textbf{Layer}            & \textbf{Output shape}   & \textbf{Params} & \textbf{MACs}  \\
\hline
0   & Input                                & $64\times96\times1$     & --      & --       \\
--  & Mean subtraction ($\lambda$)         & $64\times96\times1$     & 0       & $6\,144$ \\
1   & Conv2D(32, $3\!\times\!3$) + BN + ReLU + Pool(2)   & $32\times48\times32$     & $\sim\!320$    & $\sim\!56$\,K   \\
2   & Sep(64, $3\!\times\!3$) + BN + ReLU + Pool(2)      & $16\times24\times64$     & $\sim\!2\,400$  & $\sim\!116$\,K  \\
3   & Sep(96, $3\!\times\!3$) + BN + ReLU + Pool(2)      & $8\times12\times96$      & $\sim\!6\,720$  & $\sim\!73$\,K   \\
4   & Sep(128, $3\!\times\!3$) + BN + ReLU + Pool(2)     & $4\times6\times128$      & $\sim\!13\,152$ & $\sim\!39$\,K   \\
5   & Sep(192, $3\!\times\!3$) + BN + ReLU               & $4\times6\times192$      & $\sim\!25\,728$ & $\sim\!53$\,K   \\
--  & GlobalAveragePooling2D                & $192$                  & 0       & $4\,608$ \\
FC  & Dense(256) + ReLU + Dropout(0.35)     & $256$                  & $\sim\!49\,408$ & $49\,408$ \\
\hline
$L$ & Dense(3,  softmax)                    & $3$                    & 771     & 771     \\
$S_1$ & Dense(11, softmax)                  & $11$                   & 2\,827  & 2\,827  \\
$S_2$ & Dense(11, softmax)                  & $11$                   & 2\,827  & 2\,827  \\
$S_3$ & Dense(10, softmax)                  & $10$                   & 2\,570  & 2\,570  \\
\hline
\multicolumn{3}{r}{\textbf{Total}} & $\mathbf{\approx 109}$\textbf{\,K} & $\mathbf{\approx 408}$\textbf{\,K} \\
\hline
\end{tabular}
\end{table}

%--- Preprocessing --------------------------------------------------------
\subsubsection{Per-Image Mean Subtraction}
\label{sec:mean_sub}

Following Goodfellow~et~al.~\cite{goodfellow2014multidigitnumberrecognitionstreet}~(§4), the first
operation inside the network is per-image mean subtraction.  Let
$\mathbf{X} \in \mathbb{R}^{H \times W \times 1}$ be the input; the
normalised input is:

\begin{equation}
\label{eq:mean_sub}
  \tilde{\mathbf{X}} = \mathbf{X}
    - \frac{1}{HW}\sum_{i=1}^{H}\sum_{j=1}^{W} X_{ij}.
\end{equation}

Embedding this subtraction as a non-trainable $\lambda$-layer makes the model
\emph{self-contained}: no external preprocessing is required at inference time,
simplifying the TFLite deployment pipeline on the Raspberry~Pi~5.

%-----------------------------------------------------------------------
\subsection{Training Objective}
\label{sec:training}

Training maximises the log-probability of the correct sequence
(Eq.~\eqref{eq:prob_model}) via the cross-entropy decomposition.  Because
$L$, $S_1$, $S_2$, $S_3$ are conditionally independent given the shared
features~$\mathbf{H}$, the total loss factorises as:

\begin{equation}
\label{eq:loss}
  \mathcal{L}
  = -\log P(L \mid \mathbf{H})
    - \sum_{i=1}^{3}\log P(S_i \mid \mathbf{H})
  = \underbrace{w_L \,\mathcal{L}_L}_{\text{length head}}
    + \underbrace{w_1 \,\mathcal{L}_{S_1}
    +  w_2 \,\mathcal{L}_{S_2}
    +  w_3 \,\mathcal{L}_{S_3}}_{\text{digit heads}},
\end{equation}

where each $\mathcal{L}_{\cdot}$ is a {sparse categorical cross-entropy}
and the weights are set to $w_L{=}4$, $w_1{=}w_2{=}w_3{=}1$ to
up-weight the length prediction, since an incorrect length invariably causes a
wrong overall transcription.

Each softmax head receives the same 256-dimensional feature vector from the
shared fully connected layer but maintains its own projection weights:

\begin{equation}
\label{eq:head}
  P(S_i = c \mid \mathbf{H})
  = \frac{\exp\!\bigl(\mathbf{w}_{i,c}^{\top}\!\mathbf{H} + b_{i,c}\bigr)}
         {\displaystyle\sum_{c'} \exp\!\bigl(\mathbf{w}_{i,c'}^{\top}\!\mathbf{H}
         + b_{i,c'}\bigr)},
  \qquad c \in \mathcal{C}_i,
\end{equation}

where $\mathcal{C}_L = \{0,1,2\}$,
$\mathcal{C}_{S_1} = \mathcal{C}_{S_2} = \{0,\dots,10\}$, and
$\mathcal{C}_{S_3} = \{0,\dots,9\}$.

%-----------------------------------------------------------------------
\subsection{Log-Probability Decoding}
\label{sec:decoding}

At inference time, the predicted number is obtained by maximising the
joint log-probability over all possible lengths
(Appendix~A of~\cite{goodfellow2014multidigitnumberrecognitionstreet}):

\begin{equation}
\label{eq:decode}
  \hat{\mathbf{s}}
  = \operatorname*{argmax}_{l,\; s_1 \dots s_l}\;
    \biggl[\log P(L{=}l \mid \mathbf{X})
    + \sum_{i=1}^{l} \max_{s_i \in \{0..9\}}
      \log P(S_i{=}s_i \mid \mathbf{X})\biggr].
\end{equation}

Because of the conditional-independence assumption, the argmax over each
digit position can be computed independently, then the cumulative
log-probabilities are summed for each candidate length~$l \in \{1,2,3\}$.
The total decoding runtime is therefore $O(N)$ -- i.e.\ three additions
and one comparison -- adding negligible latency.

Concretely, for each candidate length $l$, the algorithm computes:

\begin{equation}
\label{eq:score}
  \text{score}(l) = \log P(L{=}l{-}1)
    + \sum_{i=\max(1,\,4-l)}^{3}
      \max_{s_i \in \{0..9\}} \log P(S_i{=}s_i),
\end{equation}

and selects $\hat{l} = \operatorname{argmax}_{l \in \{1,2,3\}} \text{score}(l)$.

%-----------------------------------------------------------------------
\subsection{Computational Efficiency for Edge AI}
\label{sec:efficiency}

For an ADAS system operating on edge devices like the Raspberry Pi 5, computational efficiency is just as critical as detection accuracy. To ensure real-time performance without overwhelming the hardware, the proposed multi-head CNN must be highly optimized. This section evaluates the model's efficiency by comparing its parameter count against a standard baseline, analyzing its computational cost in terms of MAC operations, and outlining the final export process to TensorFlow Lite for seamless edge deployment.

\subsubsection{Parameter Count Comparison}

The proposed model achieves the same multi-head output structure as
Goodfellow~et~al.~\cite{goodfellow2014multidigitnumberrecognitionstreet} with
$\sim\!600\times$ fewer parameters:

\begin{table}[htbp]
\centering
\caption{Comparison between the original Goodfellow~et~al.\ architecture
         and the proposed lightweight variant.}
\label{tab:compare}
\renewcommand{\arraystretch}{1.15}
\begin{tabular}{l r r}
\hline
                               & \textbf{Goodfellow et al.}  & \textbf{Proposed}   \\
\hline
Conv layers                    & 8 + 1 locally connected     & 1 Conv + 4 SepConv  \\
FC layers                      & 2 $\times$ 3\,072           & 1 $\times$ 256      \\
Output heads                   & 6  (L + S$_1$--S$_5$)       & 4  (L + S$_1$--S$_3$) \\
Total parameters               & $\sim$64\,M                 & $\sim$109\,K        \\
Max sequence length $N$        & 5                           & 3                   \\
Input size                     & $54\times 54 \times 3$      & $64\times 96 \times 1$ \\
Training infrastructure        & DistBelief (10 replicas)    & Single GPU / CPU    \\
\hline
\end{tabular}
\end{table}

\subsubsection{MAC Budget and Latency Estimate}

The total network requires approximately $408$\,K~MACs per inference.
On the Raspberry~Pi~5's Cortex-A76 cores (capable of
$\sim$4~GFLOPS at float32 via NEON), this translates to:

\begin{equation}
\label{eq:latency}
  t_{\text{inf}}
  \approx \frac{408 \times 10^{3}}{4 \times 10^{9}}
  \approx 102 \;\mu\text{s},
\end{equation}

well under the $67$\,ms budget needed for $15$\,fps real-time recognition.
Even with TFLite interpreter overhead, measured end-to-end latency stays below
$2$\,ms per frame for the float32 model.

\subsubsection{TFLite Export for the Multi-head CNN}

The trained Keras model is converted to TensorFlow~Lite using
\texttt{tf.lite.TFLiteConverter.from\_keras\_model()}.  The exported
float32 \texttt{.tflite} flatbuffer is $\sim$440\,KB, compact enough to
fit within the L2 cache of a Cortex-A76 core on the Raspberry~Pi~5.
The model is deployed via the LiteRT interpreter with \texttt{batch=1},
requiring no external preprocessing thanks to the embedded
mean-subtraction $\lambda$-layer (Eq.~\eqref{eq:mean_sub}).

%-----------------------------------------------------------------------
\subsection{Data Pipeline and Augmentation}
\label{sec:dataCNN}

The training pipeline mixes three data sources to maximise diversity while
keeping RAM usage low ($<200$\,MB peak):

\begin{enumerate}
  \item \textbf{MNIST (70\,K images):}  Individual $28{\times}28$ digits
        composited onto $64{\times}96$ canvases at randomised positions,
        simulating multi-digit layouts.
  \item \textbf{Real LED crops (23 images):}  Single-digit photographs of
        actual traffic-light seven-segment displays (red, green, yellow).
        Pre-augmented to 500 variants per digit using brightness jitter,
        gamma correction, Gaussian blur, salt-and-pepper noise, spatial
        shifts ($\pm 4$\,px), and small rotations ($\pm 10°$).
  \item \textbf{Real ROI images:}  Full composite timer-box crops from
        Vietnamese traffic cameras, labelled by filename
        (e.g.\ \texttt{37\_0004.png} $\to$ number~$37$).  Loaded from disk
        on-demand (streaming) to avoid large memory allocation.
\end{enumerate}

Synthetic composites and real ROI samples are interleaved via
\texttt{tf.data.Dataset.sample\_from\_datasets} with proportional
weighting, ensuring both domains appear in every training batch.

%-----------------------------------------------------------------------
% \subsection{Summary}

% The proposed multi-head CNN inherits the elegant probabilistic framework
% of Goodfellow~et~al.~\cite{goodfellow2014multidigitnumberrecognitionstreet} -- conditional independence
% of digit positions, log-probability decoding, per-image mean subtraction -- but
% replaces the heavyweight convolutional backbone with an ultra-light depthwise
% separable architecture totalling only $\sim$109\,K parameters and $\sim$408\,K
% MACs.  Combined with Global Average Pooling and INT8 post-training quantisation,
% the model achieves sub-millisecond inference on the Raspberry~Pi~5 CPU, making it
% suitable for real-time traffic-light timer recognition in ADAS applications
% on constrained edge devices.